\begin{exercisechunk}
\item \textbf{Many repeated tasks can hide a shared feature.}
\label[exercise]{ex:l13-task-diversity}
\exerciseprofile{2}{1}{2}{\exproof\quad\excalculation\quad\exinterpretation}{%
Distinguish the number of source tasks from coverage of the shared feature
directions.}
Let $d=3$, $k=2$, and let every source-task head have the form
$a_i=(c_i,0)^\top$, where $c_i\in\R$. Consider the two feature matrices
\[
 B=[e_1\ e_2]
 \qquad\text{and}\qquad
 \widetilde B=[e_1\ e_3].
\]
For either feature matrix, source task $i$ generates
$Y=X^\top Ba_i+\varepsilon$, where $X\sim\cN(0,I_3)$ and the same independent
noise law is used in both models.
\begin{enumerate}[(a),beginpenalty=10000]
\item Prove that the joint law of $(X,Y)$ for every source task is identical
under $B$ and $\widetilde B$.
\item For uniform sampling from $M$ source tasks, compute the task-diversity
matrix
\[
 D=\frac1M\sum_{i=1}^M a_ia_i^\top
\]
and its eigenvalues.
\item Let the new-task head be $a_{\rm new}=(0,1)^\top$. Compute its ambient
parameter under each feature matrix. Explain why no procedure using only the
source data can determine which of these two new-task parameters is correct.
\end{enumerate}
\end{exercisechunk}
